\bookchapter{제17장: 스마트폰과 다른 기기의 연동}

\chapterintro{당신 손의 스마트폰은 사실 슈퍼 리모컨 겸 데이터 센터입니다. 전화뿐 아니라, 컴퓨터, TV, 시계, 스피커, 자동차를 하나로 연결하여, 이 기기들이 더 이상 고립되지 않게 합니다. 이 장에서는 그들을 “팀 플레이”하게 만드는 법을 이야기합니다.}

\booksection{하나, 컴퓨터와 협업: 스마트폰과 컴퓨터는 더 이상 두 세계가 아님}

솔직히, 예전에 가장 짜증났던 건 스마트폰으로 사진을 찍고 컴퓨터로 옮겨 편집하려면, 데이터 케이블을 찾고 전송을 기다려야 했던 겁니다. 나중에 알았지만, 사실 케이블이 전혀 필요 없고, 좋은 연동은 케이블보다 10배 빠릅니다.

\begin{itemize}
  \item \textbf{Apple 생태계(AirDrop + 범용 클립보드)}: iPhone과 Mac을 쓴다면, AirDrop은 가장 먼저 배워야 할 기능. 스마트폰에서 공유 클릭 AirDrop Mac 선택, 사진이 순식간에 전송, 데이터나 Wi-Fi 불필요, Bluetooth와 Wi-Fi 다이렉트 사용. 또 잘 안 알려진 "범용 클립보드"——스마트폰에서 텍스트 복사, Mac에서 바로 붙여넣기, 마법 같습니다.
  \item \textbf{Android+Windows(Windows 연결)}: Microsoft의 "Windows 연결" 앱, 삼성 스마트폰은 내장, 다른 Android 스마트폰은 앱 스토어에서 다운로드. 페어링 후 컴퓨터에서 직접 스마트폰 알림, 문자, 사진 확인, 전화까지 가능. 2년째 쓰고 있는데, 회의 때 스마트폰을 계속 내려다보지 않아도, 컴퓨터 화면에 알림이 뜨면 충분.
  \item \textbf{브랜드 다중 화면 협업}: 화웨이의 "다중 화면 협업"과 샤오미의 "샤오미 묘향"은 더 강력——스마트폰 화면을 직접 컴퓨터에 미러링, 마우스로 파일을 스마트폰에서 컴퓨터로 드래그, 컴퓨터에서 Douyin(틱톡), 위챗 사용도 가능. 크로스 디바이스 복사·붙여넣기, 파일 드래그 모두 가능, 경험 매우 매끄러움.
\end{itemize}

\booksection{둘, TV로 화면 미러링: 온 가족이 함께 보기}

명절에 고향 갔을 때, 사촌 여동생이 TV에서 스마트폰 안의 졸업 여행 영상을 재생하고 싶어 했습니다. 한참 헤매다가, 사실 매우 간단하다는 것을 발견:

\begin{itemize}
  \item \textbf{iPhone TV}: 우측 상단에서 아래로 제어 센터 끌어내려, 두 사각형 겹친 아이콘(화면 미러링) 클릭, TV 이름 찾아 선택하면 연결. 전제는 TV와 스마트폰이 같은 Wi-Fi 연결, 그리고 TV가 AirPlay 지원(대부분 스마트 TV 지원).
  \item \textbf{Android TV}: 브랜드마다 부르는 이름이 다름, 화웨이는 "무선 화면 미러링", 샤오미는 "투사", OPPO/vivo는 "다중 화면 상호작용". 보통 빠른 스위치 패널에서 찾을 수 있고, 못 찾으면 설정에서 "화면 미러링" 검색.
  \item \textbf{비디오 앱 직접 화면 미러링}: iQiyi, Bilibili, YouTube를 TV에서 보고 싶다면 더 간단——재생 화면에서 Wi-Fi 기호가 있는 사각형 작은 아이콘 찾기, 클릭, TV 선택, 영상 전송. 이때 스마트폰은 다른 작업 가능, 영상은 TV에서 계속 재생, 상호 간섭 없음.
  \item \textbf{화면 미러링 실패 시}: 90\%는 같은 Wi-Fi에 연결되지 않아서. 나머지 10\%는 TV의 화면 미러링 기능이 꺼져 있어서, TV 설정에서 "다중 화면 상호작용" 또는 "AirPlay" 스위치 찾기.
\end{itemize}

\booksection{셋, 스마트워치/밴드와: 손목 위의 건강 관리자}

어머니께 밴드를 사드렸더니, 처음엔 “나 운동도 안 하는데 왜 사”라고 하시더라고요. 2주 쓰고 나서, 매일 5시간밖에 안 잔다는 걸 발견하고 스스로 일상을 조정하기 시작했습니다. 이게 시계/밴드의 가장 실용적인 가치——보이지 않는 건강 데이터를 보이는 숫자로 만드는 것.

\begin{itemize}
  \item \textbf{알림 필터}: 스마트폰이 가방이나 주머니에서 계속 진동할 때, 손목 들어 보면——사장님의 위챗인지 복숭아 할인인지 1초 판단, 스마트폰 꺼낼 가치 없는 메시지는 그냥 두기.
  \item \textbf{건강 데이터 동기화}: 시계는 24시간 심박, 혈산, 수면, 걸음 수 모니터링, 이 데이터는 자동으로 스마트폰의 건강 앱에 집계. 시간이 지나면 자신의 추세를 볼 수 있어——이번 달 너무 안 걸었나, 이번 주 수면 질 좋아졌나.
  \item \textbf{교통카드/출입카드 동기화}: 스마트폰에 바인딩한 버스카드와 출입카드, 대부분 브랜드는 자동으로 시계에 동기화. 지하철 탈 때 스마트폰 꺼낼 필요 없이, 손목을 개찰구에 대면 지나감, 정말 "미래가 왔구나"는 느낌.
  \item \textbf{스마트폰 찾기 도구}: 시계에는 보통 "스마트폰 찾기" 기능이 있어, 클릭하면 스마트폰이 울림. 제가 이 기능 쓰는 빈도는 예상보다 훨씬 높습니다——스마트폰을 아무 데나 놓고 어디 뒀는지 잊어서.
\end{itemize}

\booksection{넷, 스마트 홈과: 스마트폰은 집의 “총 스위치”}

스마트 홈은 듣기엔 복잡하지만, 핵심은 세 가지: 스마트폰으로 원격 제어, 자동화 규칙 설정, 음성 조작.

제 실제 사용 사례 몇 가지를 말씀드리면, 얼마나 실용적인지 알 겁니다:

\begin{itemize}
  \item \textbf{에어컨}: 여름 귀가 길에 미리 스마트폰으로 에어컨 켜기, 들어서면 시원함. 100위안 남짓에 에어컨 동반 플러그 구매면 가능, 에어컨 교체 불필요.
  \item \textbf{조명}: 자기 전 침대에 누워 일어나기 귀찮을 때, 스마트폰 클릭 또 "Hey Siri/샤오아이, 불 꺼줘".
  \item \textbf{로봇 청소기}: 외출 전 스마트폰에서 "청소 시작" 클릭, 돌아오면 바닥이 깨끗. 매일 정시 청소 설정도 가능.
  \item \textbf{카메라}: 집에 노인, 아이, 반려동물이 있다면, 스마트폰으로 언제든 집 상황 확인.
\end{itemize}

어떤 생태계 선택? 한마디 요약:

\begin{itemize}
  \item iPhone 사용 Apple HomeKit 지원 기기 우선(비싸지만 프라이버시 보호 우수).
  \item 기기 많고 가성비 원함 샤오미 미지아 생태계가 가장 다양, 조명, 플러그, 로봇 청소기 등 모두 있음.
  \item 화웨이 스마트폰 사용 HarmonyOS 스마트 연결 우선, 화웨이 기기와 경험 최고.
\end{itemize}

\booksection{다섯, Bluetooth 이어폰과 스피커: 무선의 자유}

유선 이어폰에서 Bluetooth 이어폰으로 바꾼 건, 제게 스마트폰 경험 업그레이드 중 가장 뚜렷한 하나였습니다. 더 이상 가방에서 엉킨 선을 꺼내지 않아도 되고, 문 손잡이에 걸릴 걱정도 없습니다.

\begin{itemize}
  \item \textbf{페어링 간단}: 이어폰 첫 전원 켜기 시 보통 자동 페어링 모드 진입(점등), 스마트폰 Bluetooth 켜면 검색 가능, 클릭 한 번에 연결. 이후 꺼낼 때마다 자동 연결, 아무것도 안 해도 됨.
  \item \textbf{같은 브랜드 시너지}: AirPods와 iPhone 페어링 시 자동 팝업으로 배터리 표시, 여러 Apple 기기 간 원활 전환. 화웨이 FreeBuds와 화웨이 스마트폰, 샤오미 이어폰과 샤오미 스마트폰도 유사한 팝업과 최적화.
  \item \textbf{Bluetooth 스피커}: 가족 모임, 야외 캠핑 시 Bluetooth 스피커 하나로 분위기 두 배. 연결 후 스마트폰이 당신의 DJ 부스.
  \item \textbf{지연 주의}: 게임할 때 Bluetooth 이어폰을 쓰면 약간의 지연(총소리 반초 후 들림)이 있을 수 있음. 지연에 민감하다면, 유선 이어폰 사용 또는 "저지연 모드" 지원 Bluetooth 이어폰 구매.
\end{itemize}

\begin{tcbinfo}{UWB 초광대역 기술——Bluetooth보다 더 정밀한 "공간 감각"}
Bluetooth 외에도, 신형 스마트폰에는 UWB(Ultra-Wideband, 초광대역) 칩이 숨어 있습니다. 그 능력은\textbf{정밀 거리 측정과 공간 위치 파악}——다른 기기가 얼마나 떨어져 있는지, 어느 방향에 있는지 계산, 정밀도는 센티미터급. Bluetooth는 “근처에 있다”만 알려주지만, UWB는 “왼쪽 앞 1.2미터”까지 알려줍니다. 일반적 응용:

\begin{itemize}
  \item \textbf{AirTag/찾기 정밀}: iPhone이 AirTag를 찾을 때 "7미터, 왼쪽으로 가라"는 화살표 내비게이션을 표시하는 것이 UWB 덕분. 일반 Bluetooth는 "범위 내에 있는지만" 알려줌.
  \item \textbf{디지털 차량 키}: UWB 탑재 iPhone과 Android 플래그십은 차량 키로 사용 가능, 차에 가까이 가면 자동 잠금 해제, 멀어지면 자동 잠금, 운전석과 조수석 구분도 가능, 릴레이 공격 방지.
  \item \textbf{실내 내비게이션}: 쇼핑몰, 공항에서 UWB 비콘으로 "통로급" 내비게이션 구현, GPS보다 훨씬 정확.
  \item \textbf{지원 모델}: iPhone 11 이후 모델은 모두 UWB 칩 탑재(U1/U2); Android 진영 샤오미, 삼성, vivo, OPPO 일부 플래그십도 탑재.
\end{itemize}

\end{tcbinfo}

\begin{tcbinfo}{LE Audio와 Auracast——Bluetooth 오디오의 신세대 표준}
20년 쓴 클래식 Bluetooth 오디오(A2DP)가\textbf{LE Audio}(저전력 오디오)로 교체되고 있습니다, 이는 Bluetooth 5.2 이후의 신세대 오디오 표준:

\begin{itemize}
  \item \textbf{더 절전, 음질 더 좋음}:완전히 새로운 LC3 인코딩 채택, 구표준 SBC보다 음질 더 좋음, 전력 소모 더 낮음, 이어폰 배터리 타임 눈에 띄게 연장.
  \item \textbf{듀얼 채널 오디오}: 좌우 이어버드 독립 전송, 단이어 모드 전환 시 끊김 없음.
  \item \textbf{Auracast 방송 오디오}: LE Audio의 시그니처 기능. 하나의 스마트폰/TV가\textbf{무수히 많은}이어폰에 동시 오디오 방송 가능——공항 대기실, 박물관, 헬스장, 교실에서 하나의 발신원으로, 여러 명이 자기 이어폰으로 같은 화면을 들을 수 있음.
  \item \textbf{현황}: 2024-2026년 신출시 플래그십 이어폰과 스마트폰은 보편적으로 LE Audio 지원, 구형 기기는 하드웨어 업그레이드 필요, 소프트웨어로 잠금 해제 불가. 새 이어폰 구매 시 LE Audio/Auracast 표시 지원 여부 확인.
\end{itemize}

\end{tcbinfo}

\booksection{여섯, 차량 연동 시스템과: 운전 중의 안전 파트너}

운전 중 스마트폰 보는 것은 위법이며, 목숨을 위협합니다. 하지만 때로는 정말 내비게이션, 전화, 음악이 필요합니다. 이때 CarPlay(Apple)와 Android Auto(Android)가 제격입니다.

\begin{itemize}
  \item \textbf{어떻게 연결}: 차량이 지원한다면(근년 대부분 신차 지원), 데이터 케이블로 스마트폰을 차량 USB 포트에 꽂으면, 차량 화면이 자동으로 CarPlay 또는 Android Auto 인터페이스로 전환. 일부 신차는 무선 연결 지원, 더 편리.
  \item \textbf{뭘 할 수 있나}: 내비게이션(차량 내비게이션보다 훨씬 좋음), 전화(운전대에서 수신), 음악 감상(스마트폰의 음악 앱을 차량으로 직접 미러링), 음성 위챗(Siri 또는 Xiaoai가 메시지 읽어주고 회신).
  \item \textbf{구형 차도 사용 가능}: 차가 오래되어 지원하지 않는다면, 수백 위안에 외장 CarPlay 화면을 사서 대시보드에 붙이면, 경험은 비슷.
\end{itemize}

\booksection{일곱, 연동 중의 보안과 프라이버시}

기기가 많이 연결될수록, 프라이버시 유출 위험면도 커집니다. 몇 가지 주의:

\begin{itemize}
  \item 공공장소에서 화면 미러링 전, 스마트폰 알림 표시줄에 사적 메시지가 없는지 먼저 확인, 온 가족이 보지 말아야 할 것을 보게 하지 마세요.
  \item 기기를 다른 사람에게 주거나 팔기 전, 먼저 바인딩 해제——스마트폰 설정에서 해당 기기 찾기, "제거" 또는 "페어링 취소" 선택.
  \item 스마트 홈 카메라는 유명 브랜드, 암호화 지원 제품 구매, 수십 위안 짜리 잡다한 브랜드는 피하기. 해커가 카메라를 실시간 방송한 뉴스가 너무 많습니다.
\end{itemize}

\begin{tcbpractice}{체험: 스마트폰 화면을 TV로 미러링하고, 스마트 씬 설정하기}
\begin{enumerate}
  \item 스마트폰과 TV가 같은 Wi-Fi에 연결되었는지 확인.
  \item TV에서 "화면 미러링" 또는 "AirPlay" 기능 찾아 켜기.
  \item 스마트폰에서: iPhone은 제어 센터에서 "화면 미러링" 클릭; Android는 하단에서 "무선 화면 미러링" 찾기.
  \item 가장 좋아하는 사진 한 장을 TV에 미러링, 가족에게 보여주기.
  \item 스마트 홈 기기가 있다면, 스마트폰 앱에서 "외출 모드" 생성 시도: 원 클릭으로 조명 끄기, 에어컨 끄기, 로봇 청소기 시작.
\end{enumerate}

\end{tcbpractice}

\begin{tcbinfo}{사례 1: 배터리 타임 최적화 실전}
\begin{itemize}
  \item \textbf{시나리오}: 사용자 샤오미 14 배터리 타임 저조, 반나절이면 방전.
  \item \textbf{단계}: (1) 설정 배터리 전력 소모 순위 확인 (2) 비정상 백그라운드 앱 식별 (3) 5G에서 4G로 전환 (4) 백그라운드 제한 (5) 절전 모드 활성화 (6) 충전 시 케이스 벗기 방열.
  \item \textbf{예상 결과}: 배터리 타임 30-50\% 연장, 배터리 건강도 데이터 가독.
\end{itemize}

\end{tcbinfo}
